\documentclass[letterpaper]{article}
    \usepackage{geometry}
    \usepackage{fullpage}
    \usepackage[T1]{fontenc}
    \usepackage[utf8]{inputenc}
    \usepackage{hyperref} % pdf metadata
    \usepackage{palatino}
    \usepackage{tagging}  % swich between cv types
    \usepackage{array}    % used by tables

    % insert QR code
    %   https://tex.stackexchange.com/questions/66231/how-to-put-a-picture-to-the-top-right-of-a-page
    \usepackage{graphicx} 
    \usepackage{eso-pic} 

    \pagestyle{empty}
    \textheight=10in
    \raggedright

    % Ragged right in tables
    %   see: https://tex.stackexchange.com/questions/12703/how-to-create-fixed-width-table-columns-with-text-raggedright-centered-raggedlef
    \newcolumntype{R}[1]{>{\raggedleft\let\newline\\\arraybackslash\hspace{0pt}}p{#1}}
    \newcolumntype{L}[1]{>{\raggedright\let\newline\\\arraybackslash}p{#1}}
    \newcolumntype{C}[1]{>{\centering\let\newline\\\arraybackslash\hspace{0pt}}p{#1}}



%%%%%%%%%%%%%%%%%%%%%%%%%%%%%%%%%%%%%%%%%%%%%%%%%%%%%%%%%%%%%%%%%%%%%
% MACROS

    \newcommand{\lineunder} {
        \vspace*{-8pt} \\
        \hspace*{-18pt} \hrulefill \\
    }

    \newcommand{\header} [1] {
        \vspace{10pt}
        {\hspace*{-18pt}\vspace*{6pt} \textsc{#1}}
        \vspace*{-6pt} \lineunder
        \vspace{3pt}
    }

    \newcommand{\cvprofile}[5]{
        \begin{center}
            {\Huge \scshape {#1}}\\
            %\vspace{4pt}
            {\large \sc {#2}}\\
            %\vspace{4pt}
            #3
            $\cdot$ #4
            $\cdot$ #5
        \end{center}
    }

    \newcommand{\employer}[4]{
        \vspace{5pt}
        \textbf{#1}  % company name
        \hfill #2\\  % city
        \textit{#3}  % title
        \hfill #4\\  % dates
        \vspace{1mm}
    }

    \newcommand{\jobintro}[1]{
        #1
        \vspace{5pt}
    }

    \newenvironment{jobtasklist}
        {
            \vspace{-12pt}
            \begin{itemize} \itemsep 0pt
        }{
            \end{itemize}
            \vspace{-3pt}
        }

    \newcommand{\university}[3]{
        \textbf{#1}  % name
        \hfill #2\\  % city
        #3\\         % degree
        \vspace{2mm}
    }

    \newenvironment{skillslist}
        {
            \begin{tabular}[t]{ l L{5.17in} }
        }{
            \end{tabular}
        }


%%%%%%%%%%%%%%%%%%%%%%%%%%%%%%%%%%%%%%%%%%%%%%%%%%%%%%%%%%%%%%%%%%%%%
% VARIABLES

    % QR CODE 
    \usetag{QRCODE}

    % CV TYPE
    %\usetag{DATACV}
    \usetag{GISCV}

    % PERSONAL DETAILS 
    %    CVname
    %    CVaddress
    %    CVemail
    %    CVphone
    \input _private.tex

    \newcommand{\CVjobTitle}{GIS Analyst}
    \newcommand{\CVcompany}{EDF Renewables}

    % PDF METADATA
    \hypersetup{
        pdfauthor={\CVname},
        pdftitle={\CVname~- \CVjobTitle},
        pdfsubject={\CVjobTitle}
    }

\hyphenation{
    }



%%%%%%%%%%%%%%%%%%%%%%%%%%%%%%%%%%%%%%%%%%%%%%%%%%%%%%%%%%%%%%%%%%%%%
% DOCUMENT

\begin{document}
%\vspace*{-20pt}
    

%================== Profile ==================%
\tagged{QRCODE}{
    \AddToShipoutPicture*
        {\put(496,687)
            {\includegraphics[width=1.7cm,height=1.7cm]
            {linkedin-qr3.png}}}
    }

\vspace*{-30pt}

\cvprofile
    {\CVname}
    {\CVjobTitle}
    {\CVaddress}
    {\CVemail}
    {\CVphone}

\vspace*{-6pt}


%================== Experience ==================%
\header{Experience}

%------------------ Freelance ------------------%
\tagged{DATACV}{ 
    \employer
        {Freelance}
        {San Diego, CA}
        {Data Analyst Consultant}
        {January 2019 - Current}
    }
\tagged{GISCV}{ 
    \employer
        {Freelance}
        {San Diego, CA}
        {GIS Analyst Consultant}
        {January 2019 - Current}
    }

\tagged{DATACV}{
    \jobintro{
        Creating, managing and analyzing targeted surveys.  
        %Working closely with clients to ensure data results and collection process are aligned
        %to their needs and technical and functional requirements.
        Analyzing survey data for trends, correlations, and data quality.
        %Analyzing results of survey data along with the success of survey campaign.
        Using Python for data munging, process management, and statistical analysis.
        Tableau dashboards built for clients to explore and visualize their survey results.
        }
    }
\tagged{GISCV}{
    \jobintro{
        Providing geospatial services, producing digital and physical maps, 
        collecting field-based data, conducting spatial data modeling and analysis. 
        %Working closely with clients to develop their requirements and spatial data needs.  
        Expert use of ArcGIS Pro, ArcGIS Desktop, and ArcGIS Online.
        }
    }


%------------------ BC Transit ------------------%
\tagged{DATACV}{ 
    \employer
        {BC Transit} 
        {Victoria, Canada}
        {Data Analyst}
        {January 2013 - October 2018}
    }
\tagged{GISCV}{ 
    \employer
        {BC Transit} 
        {Victoria, Canada}
        {GIS Analyst}
        {January 2013 - October 2018}
    }

    %\begin{taggedblock}{DATACV}
\jobintro{
    \tagged{DATACV}{
            Identified ongoing planning and administrative requirements, sourced and
            developed relevant data sets and provided statistical and spatial analysis to
            %developed relevant data sets and provided statistical analysis in Python to
            support business processes. 
            Developed ETL pipelines in Python and SQL to manage and report on data
            quality in corporate enterprise data warehouse.
            }
    \tagged{GISCV}{
            Produced maps, posters, and conducted spatial data analysis for planners and
            administrators.  
            Developed spatial data models to support transportation planning and decision making. 
            Managed updates and integration of new data into corporate spatial data warehouse.
            }
    }

\begin{jobtasklist}
    \item Received 2015 corporate Recognizing Excellence and Values award 
        %for helping to create cross-departmental team that analyzed and established
        %quality measures for data collection that significantly improved data
        %quality.

    \item Implemented reporting system dashboard to analyze
    regional transport ridership, processing 150,000+ transactions per day from
    a 700+ vehicle fleet
    
%    \tagged{DATACV}{
%        \item Developed performance metrics system in Python to monitor security and
%        quality of regional reporting systems
%    }
\end{jobtasklist}


%------------------ World Bank ------------------%
\tagged{DATACV}{ 
    \employer
        {World Bank} 
        {Washington, D.C.}
        {Consultant}
        {June 2012 - December 2012}
    }
\tagged{GISCV}{ 
    \employer
        {World Bank} 
        {Washington, D.C.}
        {GIS Consultant}
        {June 2012 - December 2012}
    }

Deployed \tagged{GISCV}{spatial} data sharing platforms to multiple countries to
aid with national disaster risk mitigation efforts. Identified and cleaned
national data sets, and ensured standardized metadata records.
\vspace{5pt}

\begin{jobtasklist}
    \item Helped government ministries planning for disaster mitigation by:
    defining data needs with project stakeholders, consolidating demographic,
    socioeconomic, and geographic data sets, and building 
    geodatabases
    
%    \item Developed data cleaning and database uploading tools in Python to 
%    standardize data cleaning and data management processes
\end{jobtasklist}


%------------------ Development Alternatives ------------------%
% \tagged{DATACV}{ 
%     \employer
%         {Development Alternatives Inc.~(DAI)} 
%         {Washington, D.C.}
%         {Principal Development Specialist}
%         {December 2006 - July 2010}
%     }
% \tagged{GISCV}{ 
%     \employer
%         {Development Alternatives Inc.~(DAI)} 
%         {Washington, D.C.}
%         {Principal GIS Development Specialist}
%         {December 2006 - July 2010}
%     }
% 
% \jobintro{
%     Developed and deployed mapping and information-sharing applications as part of
%         \tagged{DATACV}{data team.}
%         \tagged{GISCV}{GIS / mapping team.}
%     Visited field offices to create, train, and assess data collection and analysis
%     field units. 
%     Supported RFP process with prototypes and written submissions.
%     \tagged{GISCV}{
%         %Provided mapping and analysis services to staff and external partners.
%         Collected, converted, and consolidated spatial data sets in a variety of
%         formats into corporate spatial data warehouse. 
%         }
%     }
% 
% \begin{jobtasklist}
%     \item Promoted from
%         \tagged{DATACV}{Data Engineer to Principal Development Specialist}
%         \tagged{GISCV}{GIS Engineer to Principal GIS Development Specialist}
% 
% %    \item Worked on USAID's Pakistan activity tracking system (a
% %    web-based PHP application) that consolidated project information on
% %    thousands of partner and contractor activities in the region
% 
%     \item Designed and built 
%         %Python and GoogleEarth 
%         mapping system which
%     effectively visualized USAID's partnerships with corporations world wide for
%     the Global Development Alliance (GDA) project
% 
% %    \item Developed prototype national mapping and statistical research tool to
% %    collect and map statistics for Pakistan's national and regional
% %    legislatures
% 
% \end{jobtasklist}



%================== Skills ==================%
\header{Skills}
    \tagged{GISCV}{
        \vspace{-9pt} }
    \begin{taggedblock}{DATACV}
        \begin{skillslist}
%            Python: &
%                Pandas, Numpy, and Scikit-learn, Jupyter notebooks, and
%                developing ETL pipelines 
%                \\
            Data analysis: &
                Statistical analysis, cleaning, and ETL with Python and R, 
                and creating data presentations and exploratory dashboards with
                Tableau
                \\
            Databases: & 
                Oracle and PostgreSQL RDBMSs with advanced use of SQL, stored
                procedures, complex queries, and database management
                \\
%            ETL: & 
%                Developing data pipelines and processes for cleaning, verifying,
%                and managing data sets in Python and R
%                \\
            Communication: &
                Presented workshops and training seminars summarizing technical
                processes and analysis to technical and non-technical audiences
                with Visio and PowerPoint
                \\
            Customer service: &
                Engaged with clients and partners to identify data needs, plan data
                collection projects, and provide ongoing technical support
                \\
%            Collaboration: &
%                Engaged with planners, clients, and partners to identify data
%                needs, plan data collection projects, refine analysis and models
%                \\
        \end{skillslist}
    \end{taggedblock}
    \begin{taggedblock}{GISCV}
        \begin{skillslist}
            GIS:& 
                ESRI ArcGIS Desktop \& ArcGIS Pro, Spatial Analyst, 3D Analyst,
                spatial data pipelines with Safe FME, 
                Python for geoprocessing ArcPy scripts and tooboxes
                \\
%            Python:& 
%                Geoprocessing ArcPy scripts and toolbox development, Jupyter notebooks 
%                for analyzing data with Pandas, Numpy, and Scikit-learn
%                \\
%             Databases:& 
%                 Oracle spatial and PostgreSQL RDBMSs with advanced use of SQL,
%                 stored procedures, complex queries
%                 \\
            Communication: &
                Presented workshops and training seminars summarizing technical
                processes and analysis to technical and non-technical audiences
                \\
%            ETL: & 
%                Developing data pipelines and processes for cleaning, verifying,
%                and managing large data sets in Python
%                \\
            Customer service: &
                Engaged with clients, and partners to identify data needs, plan data
                collection projects, and provide ongoing technical support
                \\
%            Collaboration: &
%                Engaged with planners, clients, and partners to identify data
%                needs, plan data collection projects, refine analysis and models
%                \\
        \end{skillslist}
    \end{taggedblock}



%================== Education ==================%
\header{Education}
    \university
        {Queen's University}
        {Kingston, Canada}
        {MSc.~Computer Science}

    \university
        {University of Ottawa}
        {Ottawa, Canada}
        {Graduate Diploma International Development}

    \university
        {Queen's University}
        {Kingston, Canada}
        {BSc.~Psychology \& Computer Science}

\end{document}


\iffalse
Role

The GIS Analyst is responsible for creating sophisticated mapping and
visualization tools to help planners, modelers, and decision-makers visualize
real-world events. The GIS Analyst will prepare maps and analytical reports,
maintain and enhance an urban land information system, and design, develop,
test, maintain, and document GIS databases and applications to support
planning, research, and public safety programs. This position is ideal for a
GIS professional with a strong interest in furthering their expertise
performing specialized mapping and analysis in a regional planning environment.

 

Job Responsibilities

• Perform GIS analyses, develop GIS models and databases, design and prepare
maps, and create other GIS-related products.
• Design, develop, integrate, and maintain geographic information databases,
and related tables, forms, and reports used to support regional planning, asset
management, capital improvement projects, and public safety data and
application workflows.
• Support the ongoing development and maintenance of the SPACECORE database and
application tools; edit and code land inventory data in a production
environment.
• Prepare, design, and produce a variety of maps and other graphic
representations displaying layers and attribute data from databases; develop
and maintain data layers using GIS tools and relational databases.
• Perform complex spatial analysis; interpret detailed plans, aerial imagery,
maps and legal descriptions; code and digitize spatial and non-spatial data
into various GIS layers; perform quality control checks to ensure integrity.
• Perform GIS data maintenance, manipulation, analysis, extraction, and
generation; perform data research, investigation, and verification.
• Consult with staff to determine product and information needs and develop
effective, efficient responses to GIS-related requests.
• Develop and implement automated procedures and programs to make routine
geoprocessing workflows faster and more efficient; use object-oriented
programming and relational database technologies to enhance systems and add
functionality to new or existing applications.
• Assist with the ongoing development and maintenance of web sites, desktop,
and mobile applications; review and analyze current applications to identify
opportunities for enhancement.
• Develop visualization applications and tools to enhance effective
decision-making by staff and regional stakeholders.
• Prepare and maintain technical documentation of databases and geoprocessing
programs and procedures; prepare and manage metadata; maintain accurate and
organized records.
• Coordinate with state and local technical committees and agencies to share
GIS expertise; jointly develop GIS databases, applications, and resources.
• Participate in inter-departmental, inter-agency, and binational project teams
assembled for land use, demographic, and econometric modeling, analyses, and
planning.
• Prepare and present written, oral, and visual reports to policy and
stakeholder committees, upper management, member agencies, and working groups.

 

Experience and Qualifications

• A bachelor’s degree with major course work in geography, planning, computer
science, information technology, software engineering, or a related field A
Master’s degree is desirable. A combination of education and recent work
experience may be considered in lieu of a degree
• One to three years of recent career experience in a GIS or a
research-oriented field.
• Demonstrated knowledge of GIS concepts, principles, practices, and
techniques, including cartography, design principles, and map creation
techniques.
• Experience with ESRI’s ArcGIS for Desktop software; experience using ArcGIS
for Server, ArcGIS Pro, Spatial Analyst, 3D Analyst, and Network Analyst
    extensions are desirable.
• Experience designing geodatabases with ArcGIS architecture; experience
compiling and integrating data from multiple GIS sources.
• Demonstrated experience with relational database management systems such as
Microsoft SQL Server; a working knowledge of query design using the SQL
programing language and knowledge of Python is desirable.
• Working knowledge of GIS web services, ArcGIS Online (AGOL), and web mapping
and application development are desirable.
• Working knowledge of emerging GIS technologies for data management, analysis,
and visualization.
• Strong computer skills and proficiency with the Windows operating system and
Microsoft Office applications, especially Word, Excel, and PowerPoint.
• Demonstrated ability to maintain and effectively document GIS databases,
models, programs, and applications.
• Ability to work independently and develop GIS procedures that improve
process, workflow, and analytic capabilities.
• Knowledge of technologies and methods used to create digital orthorectified
imagery and terrain data and basic image processing concepts is desirable.
• Experience working with end-users to determine GIS needs and developing
applications, maps, and other products that best meet those needs.
• Excellent organizational skills, attention to detail, and the ability to
maintain accurate records and work on concurrent projects.
• Strong interpersonal, written, and oral communication skills, including the
ability to effectively communicate technical information to non-technical
audiences.
 
\fi


